\documentclass[24pt, aspectratio=169]{beamer}
\usepackage{pgfpages}
%\setbeameroption{show notes on second screen=right} % Both
\setbeameroption{hide notes} % Only slides
%%%%%%%% tema e cor %%%%%%%%
\mode<presentation> {
\usetheme{Boadilla}
\usecolortheme{seahorse}
%\usepackage{mathptmx}
\usepackage{relsize} % for bigger size summations
\usepackage{subcaption}
\usepackage{bm}
% \usepackage{inputenc}
% \usepackage{csquotes}
\usepackage{amsmath}
\usepackage{amssymb}
\usepackage{amsfonts}
\usepackage{mathtools}
\usepackage{xcolor}

\setbeamertemplate{footline}[frame number]{}
\setbeamertemplate{navigation symbols}{}
%\setbeamertemplate{footline}{}
\setbeamertemplate{footline}{
	\raisebox{5pt}{\makebox[\paperwidth]{\hfill\makebox[10pt]{\tiny\insertframenumber}}}}

\newcommand\unnumbered{\setbeamertemplate{footline}{}}





\usepackage[sorting=none,maxcitenames=9,maxbibnames=9,backend=bibtex,firstinits=true,style=numeric]{biblatex}
\AtBeginBibliography{\scriptsize}
\addbibresource{ref.bib}

\usepackage{booktabs}
\setlength{\heavyrulewidth}{1.5pt}
\setlength{\abovetopsep}{4pt}
\usepackage{multirow}

%color for the tavble row
\definecolor{LightCyan}{rgb}{0.88,1,1}
\definecolor{Gray}{gray}{0.9}

%\usecolortheme{albatross}
\DeclareMathOperator{\Ra}{R_1}
\DeclareMathOperator{\Rb}{R_2}
\DeclareMathOperator{\Rc}{R_3}
\DeclareMathOperator{\thh}{th}
\DeclareMathOperator{\score}{score}
\DeclareMathOperator{\iou}{IoU}
\DeclareMathOperator*{\argmax}{arg\,max}
}
\DeclareMathOperator*{\minimize}{minimize}
% Presentation configuration
\usepackage[ngerman, english]{babel}
\usepackage[utf8]{inputenc}
\usepackage{graphicx}
\usepackage{algorithm}
\usepackage{algorithmic}

\institute[TUK]{
\begin{Large}
    HyPo Conference\\
\end{Large}
\vspace{0.5cm}
}
\date{\today}
\captionsetup{margin=25pt,font=small,labelfont=bf}



%%%%%%%% titulo e subtitulo %%%%%%%%
\title[Conference Presentation]{Interpretable Multimodal Anomaly Detection for Powder-Fed Laser-Based DED}
\date{}
%%%%%%%% nome dos autores %%%%%%%%
\author[]{
} 
\titlegraphic{\includegraphics[width=4cm,height=1cm,keepaspectratio]{img/logo/mlgroup.png}
}

\begin{document}
% Title slide
{\unnumbered
\begin{frame}
\titlepage
\note[item]{Hi everyone, its another friday afternoon, the research talk session of ML group,  I would like to welcome you all here. My name is Saurabh Varshneya and I am a PhD student at this group. Today I would like to discuss about our research on Interpretable Integration of Multiple MOdalities, and hpw we plan to direct our research in the area of DeeptensorFusion}
\note[item]{I would like to thank you for giving me the opportunity to present this  progress report}
\end{frame}
}

%%%%%%%% slides %%%%%%%%
% Section 0: Introduction
\section{Module 0: Introduction}

\begin{frame}{Anomaly Detection detects anomalous patterns in data}
\end{frame}

\begin{frame}{Deep-SVDD learns a compact hypersphere of normal data representations for Anomaly Detection}
\end{frame}

\begin{frame}{Multimodal Deep SVDD combine modalities to detect anomalies}
\end{frame}

\section{Module 1: Algorithms}
\begin{frame}{Interpretability determine relevance of modalities in linear combinations and non-linear interactions}
\end{frame}

\begin{frame}{Weight parameters of fusion layer determine relevance scores}
\end{frame}

\begin{frame}{Tensor product captures non-linear interactions}
\end{frame}

\begin{frame}{Normalization scheme disentangles the interactions with linear combinations}
\end{frame}

\begin{frame}{I-MAD fuse modalities in an interpretable manner for Anomaly Detection}
\end{frame}

\section{Module 2: Experiments and Results}% Section 2: Results
\begin{frame}{DED-LB simulation provide multimodal data for Anomaly Detection}
\end{frame}

\begin{frame}{Depth-Ratio is used for deciding labels for testing}
\end{frame}

\begin{frame}{Fusing multiple modalities }
\end{frame}

%%%%%%%% agradecimentos %%%%%%%%
%\begin{frame}{Acknowledgments}
%    \large{Under Construction}
%\end{frame}
%------------------------------------------------


%------------------------------------------------
%------------------------------------------------
%%%%%%%% referencias %%%%%%%%
\nocite{*}
\begin{frame}[allowframebreaks]{ References}
\scriptsize
\printbibliography
\end{frame}


%%%%%%%% repete primeiro slide %%%%%%%%
% \begin{frame}
% \titlepage 
% \end{frame}


\end{document}
